\label{sec:conclusions}

A framework for the calculation of correlation functions in the
perturbative gradient-flow formalism using tools and techniques from
standard perturbative \abbrev{QCD} calculations has been presented. It
employs a Feynman-diagrammatic approach to the gradient-flow formalism,
based on the five-dimensional field-theoretical formulation of
\citere{Luscher:2011bx}. We implemented the corresponding \qcd\ Feynman
rules within the gradient-flow formalism as well as those of the
composite operators relevant for the quantities considered in this
paper, and automated the reduction and calculation of algebraic
expressions and loop integrals to a large extent. An important
ingredient of our setup is the reduction to master integrals with the
help of integration-by-parts identities that also involve the flow-time
variables. Except for a few cases where an analytic result was derived,
the master integrals were evaluated numerically.

To demonstrate the viability of this setup, we reproduced and improved
the three-loop result for the gluon condensate of an earlier
calculation, which was based on explicit field-theoretic calculations at
the level of Wick contractions\,\cite{Harlander:2016vzb}.

Using this setup, we then evaluated the three-loop approximations to the
quark condensate and the conversion factor from \msbar{} renormalized to
ringed quark fields, both of which are new results. Checks on gauge
parameter independence, renormalization group invariance, as well as
alternative calculational approaches (with and without reduction to
master integrals) convinced us of the correctness of these results.

The gluon and the quark condensate lend themselves for straightforward
definitions of a gradient-flow gauge coupling and a gradient-flow
quark mass, for which we derive the three-loop matching relations to the
$\msbar$ scheme. Provided that sufficiently accurate lattice results for
these quantities become available, the conversion factors evaluated in
this paper should allow for precise first-principle determinations of
the gauge coupling,\footnote{First steps in this direction have been
  taken in \citere{Lambrou:2016uet}.} and possibly also quark masses.

While focusing on the calculation of vacuum matrix elements in this
paper, our setup should be applicable in principle to more general
problems which may involve a larger number of scales, for example.
